% Part: first-order-logic
% Chapter: models-theories
% Section: expressing-relations

\documentclass[../../../include/open-logic-section]{subfiles}

\begin{document}

\olfileid[hi]{fol}{mat}{exr}

\olsection{\article{structure} \printtoken{S}{structure} में संबंध व्यक्त करना}

\begin{explain}
!!{formula}s का एक मुख्य उपयोग !!a{structure}~$\Struct M$ के गुणधर्मों और
संबंधों को उसकी भाषा~$\Lang L$---अर्थात्~$\Struct M$ की भाषा---के मूल चिह्नों
के पदों में व्यक्त करना है। इसका हमारा आशय यह है:~$\Struct M$ का !!{domain}
वस्तुओं का समुच्चय है।
~$\Struct M$ में !!{constant}s, !!{function}s और !!{predicate}s की
व्याख्या क्रमशः~$\Domain M$ की कुछ वस्तुओं,~$\Domain M$ पर फलनों और~$\Domain M$
पर संबंधों द्वारा होती है। उदाहरणतः यदि~$\Obj A^2_0$,~$\Lang L$ में है,
तो~$\Struct M$ उसे संबंध~$R = \Assign{{\Obj
  A^2_0}}{M}$ नियत करता है। तब
सूत्र~$\Atom{\Obj A^2_0}{\Obj v_1, \Obj v_2}$ उसी संबंध को निम्न अर्थ में
\emph{व्यक्त करता है}: यदि कोई चर-नियतन~$s$,~$\Obj v_1$ को $a \in \Domain{M}$
और~$\Obj v_2$ को $b \in \Domain M$ पर प्रतिचित्रित करता है, तो
\[
Rab \text{\quad तभी और केवल तभी जब\quad} \Sat{M}{\Atom{\Obj A^2_0}{\Obj v_1, \Obj v_2}}[s].
\]
ध्यान दें कि यहाँ चर-नियतनों को सम्मिलित करना आवश्यक है: हम केवल ``$Rab$
तभी और केवल तभी जब $\Sat{M}{\Atom{\Obj A^2_0}{a, b}}$'' नहीं कह सकते,
क्योंकि~$a$ और~$b$ हमारी भाषा के चिह्न नहीं, बल्कि~$\Domain{M}$ के
!!{element}s हैं।

हमारे साधन केवल परमाण्विक !!{formula}s तक सीमित नहीं हैं; तार्किक संयोजकों और
परिमाणकों से उन्हें जोड़ भी सकते हैं। इसलिए अधिक जटिल !!{formula}s के द्वारा
ऐसे अन्य संबंध परिभाषित किए जा सकते हैं जो~$\Struct M$ में सीधे निर्मित नहीं हैं।
हमारी रुचि यह जानने में है कि ऐसा कैसे किया जाता है, और विशेषतः किसी
!!a{structure} में कौन-से संबंध परिभाषित किए जा सकते हैं।
\end{explain}

\begin{defn}
मान लें कि~$!A(\Obj v_1,\dots, \Obj v_n)$,~$\Lang L$ का ऐसा !!a{formula}
है जिसमें केवल~$\Obj v_1$,\dots,~$\Obj v_n$ मुक्त आते हैं, और~$\Struct
M$,
~$\Lang L$ की कोई !!a{structure} है।~$!A(\Obj v_1,\dots, \Obj v_n)$,
संबंध~$R \subseteq \Domain M^n$ को तभी और केवल तभी \emph{व्यक्त करता है}
जब
\[
Ra_1\dots a_n \text{\quad तभी और केवल तभी जब\quad} \Sat{M}{\Atom{!A}{\Obj
    v_1,\dots, \Obj v_n}}[s]
\]
ऐसे प्रत्येक चर-नियतन~$s$ के लिए जिसके लिए $s(\Obj v_i) = a_i$ ($i = 1,
\dots, n$)।
\end{defn}

\begin{ex}
अंकगणित के मानक प्रतिरूप~$\Struct N$ में !!{formula}~$\Obj
v_1 < \Obj v_2 \lor \eq[\Obj v_1][\Obj v_2]$,~$\le$ संबंध को~$\Nat$ पर
व्यक्त करता है।
!!{formula}~$\eq[\Obj v_2][\Obj v_1']$ परवर्ती संबंध को, अर्थात् उस संबंध
$R \subseteq
\Nat^2$ को व्यक्त करता है जिसमें~$Rnm$ तब सत्य है जब~$m$,~$n$ का परवर्ती हो।
सूत्र~$\eq[\Obj v_1][\Obj v_2']$ पूर्ववर्ती संबंध को व्यक्त
करता है। !!{formula}s के इस युग्म में~$\lexists[\Obj v_3][(\eq/[\Obj v_3][\Obj 0] \land
  \eq[\Obj v_2][(\Obj v_1 + \Obj v_3)])]$ पहला और $\lexists[\Obj
  v_3][\eq[(\Obj v_1 + {\Obj v_3}')][v_2]]$ दूसरा है; दोनों~$\Obj <$ संबंध को
व्यक्त करते हैं। इसका अर्थ है कि विधेय-चिह्न~$<$ वास्तव में अंकगणित की भाषा
में अनावश्यक है; उसे परिभाषित किया जा सकता है।
\end{ex}

\begin{explain}
यह विचार केवल तब रोचक नहीं जब कुछ विशिष्ट !!{structure}s सामने हों, बल्कि हर
उस स्थिति में सामान्य रूप से महत्त्वपूर्ण है जहाँ हम किसी अभीष्ट प्रतिरूप या प्रतिरूपों का
वर्णन करने के लिए भाषा का उपयोग करते हैं, अर्थात् जब हम सिद्धांतों पर विचार
करते हैं। इन सिद्धांतों में प्रायः मूल चिह्नों के रूप में केवल कुछ
!!{predicate}s का उपयोग होता है, किंतु जिस क्षेत्र का वे वर्णन करते हैं उसमें अनेक
अन्य संबंध भी महत्त्वपूर्ण भूमिका निभाते हैं। यदि इन अन्य संबंधों को भाषा के
मूल !!{predicate}s की व्याख्या करने वाले संबंधों द्वारा व्यवस्थित रूप से
व्यक्त किया जा सके, तो कहते हैं कि हम उन्हें भाषा में \emph{परिभाषित} कर सकते
हैं।
\end{explain}

\begin{prob}
$\Lang L_A$ में निम्न संबंधों को परिभाषित करने वाले !!{formula}s को खोजिए:
\begin{enumerate}
\item $n$,~$i$ और~$j$ के बीच है;
\item $n$,~$m$ को निःशेष विभाजित करता है (अर्थात्~$m$,~$n$ का गुणज है);
\item $n$ अभाज्य संख्या है (अर्थात्~$1$ और~$n$ के अतिरिक्त कोई संख्या~$n$ को
  निःशेष विभाजित नहीं करती)।
\end{enumerate}
\end{prob}

\begin{prob}
मान लें कि सूत्र~$!A(\Obj v_1, \Obj v_2)$ संबंध~$R
\subseteq \Domain M^2$ को किसी !!a{structure}~$\Struct M$ में व्यक्त करता है।
निम्न संबंधों को व्यक्त करने वाले सूत्र खोजिए:
\begin{enumerate}
\item संबंध~$R^{-1}$, अर्थात्~$R$ का प्रतिलोम;
\item सापेक्ष गुणनफल~$R \mid R$;
\end{enumerate}
क्या आप~$R^+$, अर्थात्~$R$ के संक्रामक संवरक, को व्यक्त करने की कोई विधि खोज सकते हैं?
\end{prob}

\begin{prob}
मान लें कि~$\Lang{L}$ वह भाषा है जिसमें केवल द्विस्थानी विधेय-चिह्न~$<$ है
(निस्संदेह~$\eq$ को छोड़कर कोई अन्य !!{constant}, !!{function} या
!!{predicate} नहीं)। मान लें कि~$\Struct{N}$ वह संरचना है जिसके लिए
$\Domain{N} = \Nat$ और $\Assign{<}{N} = \Setabs{\tuple{n,m}}{n <
  m}$।
निम्न को सिद्ध कीजिए:
\begin{enumerate}
\item $\{ 0 \}$,~$\Struct{N}$ में परिभाषणीय है;
\item $\{ 1 \}$,~$\Struct{N}$ में परिभाषणीय है;
\item $\{ 2 \}$,~$\Struct{N}$ में परिभाषणीय है;
\item प्रत्येक $n \in \Nat$ के लिए समुच्चय~$\{ n \}$,~$\Struct{N}$ में
  परिभाषणीय है;
\item $\Domain{N}$ का प्रत्येक परिमित उपसमुच्चय~$\Struct{N}$ में परिभाषणीय
  है;
\item $\Domain{N}$ का प्रत्येक सहपरिमित उपसमुच्चय~$\Struct{N}$ में
  परिभाषणीय है (जहाँ $X \subseteq \Nat$ सहपरिमित तभी और केवल तभी जब
  $\Nat \setminus X$ परिमित हो)।
\end{enumerate}
\end{prob}


\end{document}
